% Part: first-order-logic
% Chapter: tableaux
% Section: identity

\documentclass[../../../include/open-logic-section]{subfiles}

\begin{document}

\olfileid{fol}{tab}{ide}

\olsection{\usetoken{P}{tableau} مع \usetoken{S}{identity}}

تتطلب !!^{tableau}s ذات ال!!{identity} قواعد استدلال إضافية. وقواعد
$\eq$ هي الآتية (حيث $t$ و$t_1$ و$t_2$ حدود مغلقة):

\begin{defish}
\AxiomC{}
\RightLabel{$\eq$}
\UnaryInfC{\sFmla{\True}{\eq[t][t]}}
\DisplayProof
\hfill
\AxiomC{\sFmla{\True}{\eq[t_1][t_2]}}
\noLine
\UnaryInfC{\sFmla{\True}{!A(t_1)}}
\RightLabel{$\TRule{\True}{\eq}$}
\UnaryInfC{\sFmla{\True}{!A(t_2)}}
\DisplayProof
\hfill
\AxiomC{\sFmla{\True}{\eq[t_1][t_2]}}
\noLine
\UnaryInfC{\sFmla{\False}{!A(t_1)}}
\RightLabel{$\TRule{\False}{\eq}$}
\UnaryInfC{\sFmla{\False}{!A(t_2)}}
\DisplayProof
\end{defish}
لاحظ أن $\TRule{\True}{\eq}$ و$\TRule{\False}{\eq}$، بخلاف جميع القواعد
الأخرى، تتطلبان أن تكون !!{formula}s موقّعتان \emph{اثنتان} ظاهرتين بالفعل
على الفرع، وهما $\sFmla{\True}{\eq[t_1][t_2]}$ و
$\sFmla{S}{!A(t_1)}$.

\begin{ex}
إذا كان $s$ و$t$ حدين مغلقين، فإن $\eq[s][t], !A(s)
\Proves !A(t)$:
\begin{oltableau}
  [\sFmla{\False}{\formula{A}(t)}, just = \TAss
    [\sFmla{\True}{\eq[s][t]}, just = \TAss
      [\sFmla{\True}{\formula{A}(s)}, just = \TAss
        [\sFmla{\True}{\formula{A}(t)}, just={\TRule{\True}{\eq}[2, 3]}, close]
      ]
    ]
  ]
\end{oltableau}
قد يكون هذا مألوفًا باسم مبدأ إحلال المتطابقات أو قانون لايبنتس.

تثبت !!^{tableau}s أن $\eq$ متناظرة، أي إن $\eq[s_1][s_2]
\Proves \eq[s_2][s_1]$:
\begin{oltableau}
  [\sFmla{\False}{\eq[s_2][s_1]}, just = \TAss
    [\sFmla{\True}{\eq[s_1][s_2]}, just = \TAss
      [\sFmla{\True}{\eq[s_1][s_1]}, just = {$\eq$}
        [\sFmla{\True}{\eq[s_2][s_1]}, just = {\TRule{\True}{\eq}[2, 3]}, close]
      ]
    ]
  ]
\end{oltableau}
هنا السطر $2$ هو ال!!{formula} المتطلبة الأولى
~$\sFmla{\True}{\eq[s_1][s_2]}$ لــ$\TRule{\True}{\eq}$. والسطر~$3$
هو المتطلب الثاني، على الصورة $\sFmla{\True}{!A(s_1)}$---تصور أن
$!A(x)$ هي $\eq[x][s_1]$؛ عندئذ تكون $!A(s_1)$ هي $\eq[s_1][s_1]$
وتكون $!A(s_2)$ هي $\eq[s_2][s_1]$.

وتثبت أيضًا أن $\eq$ متعدية، أي إن $\eq[s_1][s_2],
\eq[s_2][s_3] \Proves \eq[s_1][s_3]$:
\begin{oltableau}
  [\sFmla{\False}{\eq[s_1][s_3]}, just = \TAss
    [\sFmla{\True}{\eq[s_1][s_2]}, just = \TAss
      [\sFmla{\True}{\eq[s_2][s_3]}, just = \TAss
        [\sFmla{\True}{\eq[s_1][s_3]}, just = {\TRule{\True}{\eq}[3, 2]}, close]
      ]
    ]
  ]
\end{oltableau}
في هذا ال!!{tableau}، ال!!{formula} المتطلبة الأولى لــ
$\TRule{\True}{\eq}$ هي السطر~$3$، $\sFmla{\True}{\eq[s_2][s_3]}$
(حيث يؤدي $s_2$ دور~$t_1$ ويؤدي $s_3$ دور~$t_2$). والمتطلب الثاني،
على الصورة~$\sFmla{\True}{!A(s_2)}$، هو السطر~$2$. وهنا تصور أن
$!A(x)$ هي $\eq[s_1][x]$؛ فيجعل ذلك $!A(s_2)$ هي
$\eq[s_1][s_2]$ (أي السطر~$2$)، ويجعل $!A(s_3)$ ال!!{formula}
~$\eq[s_1][s_3]$ في النتيجة.
\end{ex}

\begin{prob}
أنشئ !!{tableau}s مغلقة لما يأتي:
\begin{enumerate}
\item $\sFmla{\False}{\lforall[x][\lforall[y][((x = y \land !A(x))
      \lif !A(y))]]}$
\item $\sFmla{\False}{\lexists[x][(!A(x) \land
    \lforall[y][(!A(y) \lif y = x)])]}$,\\
  $\sFmla{\True}{\lexists[x][!A(x)] \land
    \lforall[y][\lforall[z][((!A(y) \land !A(z)) \lif y = z)]]}$
\end{enumerate}
\end{prob}

\end{document}
